\chapter{Vectores de Prueba}
\label{ap:vectores}

Los vectores de prueba son entradas concretas con la salida \emph{exacta} del código de producción. Los genera
\texttt{vectores.php}, un programa que va con el código privado: carga el motor tal como está (las clases del
emparejador, los factores y los tiempos por tramo, y las funciones de calibración del script de aprendizaje, sin ejecutar su programa principal) y
escribe \texttt{referencia/vectores.json}. Una implementación en cualquier lenguaje puede darse por correcta cuando
reproduce todos los vectores; \texttt{referencia/comprobar.py} lo hace para la implementación de referencia del
Apéndice~\ref{ap:referencia}.

Los valores reales se dan con todos sus decimales para comprobar también los cálculos en coma flotante; los
resultados que el método redondea deben coincidir exactamente.

\section{Funciones Básicas}

\begin{table}[ht]
\centering
\caption{Mediana Baja y Percentiles}\label{tab:v-percentil}
\begin{tabular}{@{}llr@{}}
\toprule
Valores & $p$ & Resultado\\
\midrule
5, 1, 4, 2, 3 & 0,5 & 3\\
1, 2, 3, 4 & 0,5 & 2 \quad(no 2,5: mediana baja)\\
10, 20, \dots, 100 & 0,9 & 90\\
\bottomrule
\end{tabular}
\end{table}

\begin{table}[ht]
\centering
\caption{Franja de una Fecha y un Instante}\label{tab:v-franja}
\begin{tabular}{@{}llrl@{}}
\toprule
Fecha & Día & Segundos & Franja\\
\midrule
2026-10-28 & miércoles & 28\,950 (8:02:30) & \texttt{laborable.punta\_manana}\\
2026-10-28 & miércoles & 21\,599 (5:59:59) & \texttt{laborable.madrugada}\\
2026-10-28 & miércoles & 21\,600 (6:00:00) & \texttt{laborable.punta\_manana}\\
2026-10-31 & sábado & 75\,600 (21:00:00) & \texttt{fin\_semana.noche}\\
2026-10-28 & miércoles & 90\,000 (25:00:00) & \texttt{laborable.madrugada}\\
2026-11-01 & domingo & 46\,800 (13:00:00) & \texttt{fin\_semana.mediodia}\\
\bottomrule
\end{tabular}
\end{table}

\begin{table}[ht]
\centering
\caption{Factor de Ritmo}\label{tab:v-k}
\begin{tabular}{@{}lrrl@{}}
\toprule
Línea & Segundos & $k$ & Por qué\\
\midrule
3414 & 28\,950 & 1,2 & Factor de línea 1,35, acotado\\
9999 & 28\,950 & 0,9 & Sin factor de línea: el global de la franja\\
9999 & 79\,200 & 0,6 & Global 0,4 de la noche, acotado\\
9999 & 50\,400 & 1,0 & Sin factor para el mediodía\\
03414 & 28\,950 & 1,2 & Se lee como la 3414\\
\bottomrule
\end{tabular}
\nota{Con $k_{\texttt{laborable.punta\_manana}}=0{,}9$, $k_{\texttt{laborable.noche}}=0{,}4$ y $k_{3414,\texttt{laborable.punta\_manana}}=1{,}35$, el miércoles 2026-10-28.}
\end{table}

\begin{table}[ht]
\centering
\caption{Proyección Sobre un Segmento de Prueba}\label{tab:v-proyeccion}
\begin{tabular}{@{}>{\raggedright\arraybackslash}p{0.22\textwidth}rrr@{}}
\toprule
Punto $(\varphi,\lambda)$ & $f$ & $d_\perp$ (m) & $\|B-A\|$ (m)\\
\midrule
$(43{,}2601;\,-2{,}925)$ & 0,5000000000000222 & 11,132000000369544 & 810,6880955946056\\
$(43{,}2605;\,-2{,}935)$ & 0 & 409,14505231598264 & 810,6827696075283\\
Segmento degenerado $A=B$, punto en $A$ & 0 & 0 & 0\\
\bottomrule
\end{tabular}
\nota{Segmento de $A=(43{,}26;\,-2{,}93)$ a $B=(43{,}26;\,-2{,}92)$.}
\end{table}

\noindent La longitud del segmento cambia ligeramente entre las dos primeras filas porque el factor
$\cos\varphi$ se toma en la latitud del autobús, no en la del segmento.

\begin{table}[ht]
\centering
\caption{Cubos de Longitud}\label{tab:v-cubo}
\begin{tabular}{@{}rr@{\qquad}rr@{\qquad}rr@{}}
\toprule
Metros & Cubo & Metros & Cubo & Metros & Cubo\\
\midrule
0 & 0 & 999,9 & 2 & 2\,999,9 & 4\\
299,9 & 0 & 1\,000 & 3 & 3\,000 & 5\\
300 & 1 & 1\,499,9 & 3 & 12\,000 & 5\\
& & 1\,500 & 4 & &\\
\bottomrule
\end{tabular}
\end{table}

\begin{table}[ht]
\centering
\caption{Respaldo por Longitud}\label{tab:v-respaldo}
\begin{tabular}{@{}rrlr>{\raggedright\arraybackslash}p{0.19\textwidth}@{}}
\toprule
$h$ (s) & $L$ (m) & Franja & Resultado (s) & Por qué\\
\midrule
120 & 800 & \texttt{laborable.punta\_manana} & 97,32000000000001 & $0{,}811\cdot120$\\
120 & 800 & \texttt{fin\_semana.noche} & 92,28 & $0{,}769\cdot120$ (\texttt{all})\\
300 & 2\,000 & \texttt{laborable.manana} & 450 & $2{,}0$ acotado a $1{,}5$\\
300 & 5\,000 & \texttt{laborable.manana} & 90 & $0{,}1$ acotado a $0{,}3$\\
300 & 100 & \texttt{laborable.manana} & nulo & Cubo 0 sin proporción\\
300 & desconocida & \texttt{laborable.manana} & nulo & Sin longitud\\
\bottomrule
\end{tabular}
\nota{Con las proporciones $r_{2,\texttt{all}}=0{,}769$, $r_{2,\texttt{laborable.punta\_manana}}=0{,}811$, $r_{4,\texttt{all}}=2{,}0$ y $r_{5,\texttt{all}}=0{,}1$.}
\end{table}

\section{Calibración}

\subsection{Tiempos por Tramo} Con las quince muestras siguientes (pista, franja, tramo, tiempo real), la calibración
da exactamente
\texttt{A>B}: $\{\texttt{n}:8,\ \texttt{all}:61,\ \texttt{laborable.punta\_manana}:69,\ \texttt{laborable.mediodia}:60\}$,
\texttt{C>D}: $\{\texttt{n}:3,\ \texttt{laborable.noche}:200\}$ y
\texttt{D>E}: $\{\texttt{n}:2,\ \texttt{laborable.manana}:60\}$; el tramo \texttt{B>C}, con una sola pista, no se
guarda.

\begin{center}
\small
\begin{tabular}{@{}llll@{\qquad}llll@{}}
\toprule
Pista & Franja & Tramo & $y$ & Pista & Franja & Tramo & $y$\\
\midrule
b1 & punta\_manana & A>B & 90 & m5 & mediodia & A>B & 61\\
b2 & punta\_manana & A>B & 100 & b1 & punta\_manana & B>C & 75\\
m0 & mediodia & A>B & 60 & c0 & noche & C>D & 200\\
m1 & mediodia & A>B & 58 & c1 & noche & C>D & 210\\
m2 & mediodia & A>B & 62 & c2 & noche & C>D & 190\\
m3 & mediodia & A>B & 64 & d1 & manana & D>E & 50\\
m4 & mediodia & A>B & 59 & d1 & manana & D>E & 70\\
 & & & & d2 & manana & D>E & 60\\
\bottomrule
\end{tabular}
\end{center}

\noindent El cálculo de \texttt{A>B} está en el ejemplo~\ref{ej:contraccion}. En \texttt{D>E}, la pista d1 tiene dos
muestras y cuenta la última (70); la mediana baja de $\{70,60\}$ es 60.

\subsection{Factor de Ritmo} Con una muestra por pista, de $H=600$~s y $Y=600\,\rho_j$:

\begin{center}
\small
\begin{tabular}{@{}lr@{}}
\toprule
Proporciones por pista $\rho_j$ & $k$\\
\midrule
0,80, 0,82, 0,84, 0,86, 0,88, 0,90, 0,92, 0,94 & 0,938\\
cuarenta veces 0,85 & 0,88\\
1,6, 1,7, \dots, 2,5 & 1,2\\
\bottomrule
\end{tabular}
\end{center}

\subsection{Paso Real} Con la pista $(28\,900,\,28\,880)$, $(28\,960,\,28\,940)$, $(29\,020,\,29\,030)$: para
$a=28\,950$, $P=28\,966{,}666666666668$; para $a=28\,800$, nulo.

\subsection{Aclarado} De 100 muestras de una pista se conservan las de índices 0, 3, 6, 10, 13, 16, 20, 23, 26, 30,
33, 36, 40, 43, 46, 50, 53, 56, 60, 63, 66, 70, 73, 76, 80, 83, 86, 90, 93 y 96.

\section{Corrección por Viaje}

\subsection{Consulta} La tabla de los vectores tiene el viaje \texttt{9999|S0|S4|28800|5} con $[20,40,60]$ en
\texttt{laborable} y $[-30,0,30]$ en \texttt{fin\_semana}; la línea \texttt{L2}, con 15~s en
\texttt{laborable.punta\_manana}; y por franja, 25~s en \texttt{laborable.punta\_manana} y $-10$~s en
\texttt{fin\_semana.noche}. Los tercios~\eqref{eq:tercio} de las posiciones 0 a 4 de un recorrido de cinco paradas son
0, 0, 1, 1 y 2; los de las posiciones 1 y 2 de tres, 1 y 2; la posición 9 de cinco da 2, y un recorrido sin paradas,
0. En un festivo, la capa 3 calcula con el domingo de esa semana: el lunes 12 de octubre y el sábado 17 dan el 18; el
domingo 1 de noviembre, el mismo día; y el martes 8 de diciembre, el 13.

\begin{center}
\small
\begin{tabular}{@{}lllrl@{}}
\toprule
Viaje & Fecha & Salida, posición & Resultado (s) & Por qué\\
\midrule
\texttt{9999|S0|S4|28800|5} & 2026-10-28 & 28\,800, 4 de 5 & 60 & Tercer tercio, laborable\\
\texttt{9999|S0|S4|28800|5} & 2026-10-28 & 28\,800, 2 de 5 & 40 & Segundo tercio\\
\texttt{9999|S0|S4|28800|5} & 2026-10-31 & 28\,800, 1 de 5 & $-30$ & Sábado: fin de semana\\
\texttt{L2|A|B|28800|3} & 2026-10-28 & 28\,800, 1 de 3 & 15 & Línea en su franja\\
\texttt{L2|A|B|72000|3} & 2026-10-31 & 72\,000, 1 de 3 & $-10$ & La línea no tiene esa franja\\
\texttt{L3|A|B|28800|3} & 2026-10-28 & 28\,800, 1 de 3 & 25 & Sin línea: la franja\\
\texttt{L3|A|B|50400|3} & 2026-10-28 & 50\,400, 1 de 3 & 0 & Nada para esa franja\\
\texttt{09999|S0|S4|28800|5} & 2026-10-28 & 28\,800, 4 de 5 & 60 & Se lee como la 9999\\
\bottomrule
\end{tabular}
\end{center}

\subsection{Calibración} Con los nueve desvíos siguientes (fecha, $D'$, viaje, posición y desvío en segundos), la
calibración~\eqref{eq:correccion-tablas} da para el viaje \texttt{L1|A|D|28800|4} en laborable $[127,150,420]$; para
\texttt{L1|A|D|32400|4}, $[-60,-60,-60]$ en laborable y $[30,30,40]$ en fin de semana; y para
\texttt{L2|X|Y|61200|6}, $[300,300,300]$. Las líneas quedan en \texttt{L1}: $\{$\texttt{laborable.punta\_manana}: 140,
\texttt{laborable.manana}: $-60$, \texttt{fin\_semana.manana}: 30$\}$ y \texttt{L2}:
$\{$\texttt{laborable.punta\_tarde}: 300$\}$, con los mismos valores por franja.

\begin{center}
\small
\begin{tabular}{@{}llllr@{}}
\toprule
Fecha & $D'$ & Viaje & Posición & Desvío\\
\midrule
2026-10-26 & 2026-10-26 & \texttt{L1|A|D|28800|4} & 1, 2 y 3 de 4 & 100, 160 y 220\\
2026-10-27 & 2026-10-27 & \texttt{L1|A|D|28800|4} & 1 y 3 de 4 & 140 y 2\,000\\
2026-10-26 & 2026-10-26 & \texttt{L1|A|D|32400|4} & 2 de 4 & $-60$\\
2026-10-31 & 2026-10-31 & \texttt{L1|A|D|32400|4} & 1 de 4 & 30\\
2026-10-12 & 2026-10-18 & \texttt{L1|A|D|32400|4} & 3 de 4 & 50\\
2026-10-26 & 2026-10-26 & \texttt{L2|X|Y|61200|6} & 5 de 6 & 300\\
\bottomrule
\end{tabular}
\end{center}

\noindent Para \texttt{L1} en \texttt{laborable.punta\_manana}, el desvío de 2\,000~s se recorta a 900. Ordenados,
$100,140,160,220,900$, el percentil 40 es el de la posición $\lfloor0{,}4\cdot4\rfloor=1$: $g=140$, y la línea, con el
peso 20 de la franja, también vale 140. El primer tercio del viaje (posición 1 de 4) tiene un desvío cada día, 100 y
140: $\redondeo\big((100+140+140)/3\big)=127$. El tercer tercio, con 220 y 900, da $\redondeo(1\,260/3)=420$. El
festivo del 12 de octubre cuenta como fin de semana: con el 31 de octubre, la franja \texttt{fin\_semana.manana} tiene
30 y 50, y su percentil 40 es 30; el primer tercio da $(30+30)/2=30$, el segundo toma el valor de la línea y el tercero
da $(50+30)/2=40$.

\section{El Emparejador Completo}

El viaje sintético tiene cinco paradas, \texttt{S0} a \texttt{S4}, en línea recta hacia el este sobre el paralelo
$43{,}26^\circ$ desde la longitud $-2{,}93^\circ$, a 0, 400, 1\,000, 2\,500 y 3\,200~m (longitudes
$-2{,}93$, $-2{,}925065928003549$, $-2{,}917664820008872$, $-2{,}8991620500221806$ y $-2{,}890527424028391$), con
llegada igual a salida en cada parada: $28\,800$, $28\,890$, $29\,010$, $29\,310$ y $29\,430$~s. El trazado, cuando se
usa, son las cinco paradas con sus distancias. La fecha es el miércoles 2026-10-28, la línea \texttt{9999} y el
vehículo \texttt{V1}. Los parámetros son:

\begin{itemize}
  \item tramos aprendidos: \texttt{S0>S1} $\{\texttt{n}:10,\ \texttt{all}:100,\ \texttt{laborable.punta\_manana}:110\}$
        y \texttt{S1>S2} $\{\texttt{n}:6,\ \texttt{all}:150,\ \texttt{laborable.punta\_manana}:140\}$;
  \item proporciones: cubo 2, $\{\texttt{all}:0{,}769,\ \texttt{laborable.punta\_manana}:0{,}811\}$; cubo 4,
        $\{\texttt{all}:0{,}64,\ \texttt{laborable.punta\_manana}:0{,}668\}$;
  \item factor global $k_{\texttt{laborable.punta\_manana}}=0{,}9$, sin factores de línea.
\end{itemize}

\noindent En el escenario de «horario muy lento», las horas de \texttt{S3} y \texttt{S4} pasan a $30\,810$ y
$30\,930$, y la tabla de tramos se sustituye por \texttt{S0>S1}, \texttt{S1>S2}, \texttt{S2>S3} y \texttt{S3>S4} con 110,
140, 240 y 100~s en \texttt{laborable.punta\_manana} (y \texttt{n}$=9$).

Las entradas del feed de cada escenario son:

\begin{center}
\small
\begin{tabularx}{\textwidth}{@{}l>{\raggedright\arraybackslash}X@{}}
\toprule
Escenario & Entrada\\
\midrule
1 y 2 & $\varphi=43{,}2601$, $\lambda=-2{,}9213653740062107$ (700~m), $\tgps=t_{\mathrm{rec}}=28\,950$, $n=3$\\
3 & $\varphi=43{,}260538986704994$ (60~m al norte), mismo $\lambda$ y horas, $n=3$\\
4 & sin posición, $t_{\mathrm{rec}}=28\,950$, $n=3$\\
5 & $\varphi=43{,}2601$, $\lambda=-2{,}895461496024842$ (2\,800~m), $\tgps=t_{\mathrm{rec}}=29\,400$, $n=3$\\
6, 7, 8, 10 y 11 & ninguna\\
9 & como el 1, con $\tgps=t_{\mathrm{rec}}=28\,300$\\
12 & $\varphi=43{,}2601$ (11~m al norte de \texttt{S0}), $\lambda=-2{,}93$, $\tgps=t_{\mathrm{rec}}=28\,200$, $n=2$, $\tnow=28\,210$\\
13 & $\varphi=43{,}26538986704995$ (600~m al norte), mismo $\lambda$ y horas, $n=2$\\
14 y 16 & como el 1, con $\tnow=29\,100$ y $29\,320$\\
15 & $\varphi=43{,}2601$, $\lambda=-2{,}917048061009316$ (1\,050~m, 50~m pasada \texttt{S2}), $\tgps=t_{\mathrm{rec}}=29\,020$, $n=3$, $\tnow=29\,040$\\
\bottomrule
\end{tabularx}
\end{center}

\noindent Los escenarios 10 y 11 usan la tabla de corrección de la sección anterior; el 11, además, con la fecha
marcada como festivo (un calendario de solo domingo que circula ese día por excepción). Los resultados están en la
Tabla~\ref{tab:escenarios}. Las posiciones calculadas son $\Tpos=28\,950$ con
$d_\perp=11{,}132000000369544$~m en los escenarios 1 (método «trazado») y 2 (método «tramo recto»); en el 12,
$\Tpos=\tgps=28\,800$ con la misma distancia y el método «cabecera»; en el 14 y el 16, la del 1; en el 15,
$\Tpos=29\,020$ con el método «trazado»; no hay posición en los demás.

\section{Del Feed a la Pantalla}

Las reglas de esta sección rodean al emparejador: cómo se lee el feed, a qué viaje del horario se asigna cada bus, qué
pasa con un bus que falta en una lectura o con un viaje que pasa de medianoche, qué se guarda de cada paso, qué retraso
se muestra y cómo se elige cada noche el percentil de la corrección por viaje. Sus vectores salen del mismo motor y
\texttt{comprobar.py} los comprueba igual que los anteriores.

\subsection{Lectura del Feed} Un fragmento SIRI con cuatro actividades, leído con fecha 2026-10-28. La referencia
\texttt{trp\_A0651\_737\_ABC1\_28500} da la línea \texttt{651} (sin el cero), el viaje \texttt{737} y la salida
$28\,500$; su \texttt{Delay} \texttt{PT2M30S} son 150~s; la hora de la posición, \texttt{07:02:30Z}, es
$\tgps=28\,950$ en Madrid, y la de la lectura, $t_{\mathrm{rec}}=28\,960$. La segunda actividad no trae posición, su
orden sale de \texttt{VisitNumber} y su \texttt{Delay} \texttt{-PT45S} son $-45$~s. La tercera se registró el día
anterior: entra en el feed, pero sin horas ni posición. La cuarta tiene una referencia de otro formato y no entra.

\subsection{Emparejado} Con la referencia exacta \texttt{trp\_A9999\_101\_X\_28800} del viaje \texttt{J} en la tabla de
alias:

\begin{center}
\small
\begin{tabularx}{\textwidth}{@{}>{\raggedright\arraybackslash}X l l@{}}
\toprule
Buses de la línea \texttt{9999}, viaje \texttt{101} & Viaje pedido & Bus elegido\\
\midrule
Uno con la referencia exacta y otro con otra referencia y la misma salida & \texttt{J} & el exacto\\
Sin referencia exacta, salida a 120~s & \texttt{J} & ese\\
Sin referencia exacta, salida a 180~s justos & \texttt{J} & ese\\
Sin referencia exacta, salida a 200~s & \texttt{J} & ninguno\\
Dos sin referencia exacta, a 150 y a 60~s & \texttt{J} & el de 60~s\\
Solo el de la referencia exacta de \texttt{J} & \texttt{K} & ninguno (no se reutiliza)\\
El de \texttt{J} y otro sin referencia exacta & \texttt{K} & el otro\\
\bottomrule
\end{tabularx}
\end{center}

\subsection{Buses que Faltan} Con $\tnow=29\,010$, de la lectura anterior se conservan, marcados como retenidos, el
bus visto hace 110~s y el que solo tiene hora de lectura de hace 180~s justos. No se conservan: el de hace 190~s ni el
de hace 210~s; el de un viaje que en la lectura nueva lleva otro bus; el que aparece en la lectura nueva en otro viaje;
el que no tiene hora; el de hora posterior a ahora; ni el que no tiene identificador de vehículo.

\subsection{Viajes que Pasan de Medianoche} A las 00:30 ($\tnow=1\,800$), el viaje que salió a las 23:30
($84\,600$) y el de las 24:10 del horario de ayer ($87\,000$) son de ayer; uno que sale a las 12:30 ($45\,000$, justo
12~h por delante) es de hoy, y uno a $45\,001$, de ayer. Para los de ayer, sus horas pasan a contarse desde la
medianoche de ayer: un dato de las 00:28:20 ($1\,700$) pasa a $88\,100$. Un bus sin hora de hoy no cuenta en ningún día.

\subsection{Registro de Pasos} De cada referencia se guarda «orden:hora», con la hora de la lectura o, sin ella, la de
la posición: \texttt{3:29000} y \texttt{2:28900}. No se guardan las entradas sin orden, sin referencia o sin hora. Al
leer, por viaje del horario queda el orden mayor de sus referencias (\texttt{3:29000} y \texttt{5:29100} dan orden 5 a
las $29\,100$); los valores nulos o mal formados se ignoran. Con ese paso, un viaje sin bus en el feed tiene pasadas las
paradas de orden menor:

\begin{center}
\small
\begin{tabular}{@{}lrlll@{}}
\toprule
Paso registrado & $\tnow$ & \texttt{S0}, \texttt{S1}, \texttt{S2} & \texttt{S3} & \texttt{S4}\\
\midrule
orden 4 a las $29\,300$ & $29\,400$ & pasadas, a las $29\,300$ & horario & capa 3: $29\,408$\\
orden 2 a las $29\,000$ & $28\,990$ & \texttt{S0} pasada a las $28\,990$ (ahora) & capa 3: $29\,310$ & capa 3: $29\,408$\\
\bottomrule
\end{tabular}
\end{center}

\noindent En el segundo caso, \texttt{S1} queda en el horario ($28\,890$) y \texttt{S2} con la capa 3 ($29\,110$).

\subsection{El Retraso Mostrado} Los minutos son los segundos entre 60 redondeados con mitades hacia fuera del cero:
29~s dan 0; 30, 1; 89, 1; 90, 2; $-30$, $-1$; $2\,700$, 45; y un retraso de más de $2\,700$~s no se muestra. Del bus
del escenario 1, con GPS, se muestra el retraso que llevará en su próxima parada, 10~s, aunque su retraso en la
posición sea 0; sin GPS (escenario 4), el de la parada registrada, 10~s; y con un dato de más de 10~min y sin hora de
lectura, 0.

\subsection{Elección del Percentil} Con un viaje de la línea \texttt{L1} y cinco desvíos en cada uno de cuatro días,
se aprende con los tres primeros y se prueba en el último:

\begin{center}
\small
\begin{tabular}{@{}lrrrrr@{}}
\toprule
Percentil & 0,30 & 0,35 & 0,40 & 0,45 & 0,50\\
\midrule
Error medio (s) & 144,4 & 144,4 & 147,4 & 149,2 & 151,2\\
A $\pm2$~min (\%) & 20 & 20 & 40 & 60 & 60\\
\bottomrule
\end{tabular}
\end{center}

\noindent Se elige 0,45: el que más acierta a $\pm2$~min y, en el empate con 0,50, el de menos error medio, aunque 0,30
tenga el menor. Sin corregir, el error medio es 126~s y el acierto, 60\,\%: la tabla nueva no bajaría el error medio y
no se escribe.
